\clearpage
\properlane{rightly-ordered-allegiance}{Receiving the world without worshipping its powers}
An earthly good can be real without being ultimate. This resolves the apparent collision between Cyrus's vocation and God's sole lordship, between Caesar's dues and God's claim. The whole formulary teaches a person to receive created goods, discharge real obligations and seek protection, while refusing to make any of these goods the source or final measure of life. Jerome distinguishes the ruler's liberating office from the true Christ; Chrysostom distinguishes rightful tribute from impious demands. Augustine's universal song and Aquinas's image of God extend their consequences into the life of the worshipping body.

\subsection{The instrument does not become the saviour}
The difficulty in Isaiah is not simply that God works beyond Israel's borders. The king receives a title of sacred resonance while the oracle twice identifies his ignorance. If usefulness were enough to establish holiness, the second assertion would make little sense. Jerome's insistence on the actual Persian king lets both assertions stand. A ruler can serve a divine purpose without knowing it and without acquiring divine standing himself.

This changes how benefit is received. Israel's restoration need not be denied because its human instrument is imperfect. Yet gratitude for restoration cannot turn into a judgment that everything about that instrument is good. The two mistakes are mirror images: refusing a real gift because its agent is not holy, or granting unlimited confidence to an agent because some good has come through him. The prophecy permits neither inference. It puts the act within a purpose named by God, for Jacob and for the knowledge of the nations.\footnote{Jerome, \work{Commentary on Isaiah}, XII, on 45:1--7, pp.127--129; Aquinas, \work{Super Isaiam}, chapter 45, paragraph 85005.}

Jerome's typology carries this distinction into the Christian reading. Cyrus can prefigure liberation without becoming identical with the liberator who reconciles the world. The contrast between type and fulfilment prevents a religious vocabulary from flattening the difference between an office and the Son. The Church's confession of Christ does not require the past to contain a series of lesser Christs whose every action is sacred. It recognizes genuine anticipations whose limits make the fulfilment more, not less, intelligible.

\subsection{The coin establishes a debt, not a deity}
The Gospel makes the same question tangible through a coin. The questioners hope to make obedience to God and the payment of tribute mutually exclusive. Jesus' answer refuses the forced choice. Chrysostom's emphasis on ``render'' makes the civil obligation concrete: a return is due within the common life in which the currency is used. It is inadequate to spiritualize the exchange before allowing the coin to remain money and the tax to remain a real demand.

But the answer's second clause prevents the first from becoming absolute. Chrysostom insists that a demand contrary to godliness has exceeded what can rightly be rendered to Caesar. The ruler's office identifies the claimant; it does not automatically establish the moral character of every claim. This is why the answer is more demanding than a convenient division into a secular half and a sacred half. The person doing the rendering belongs to God while acting within civil life.\footnote{Chrysostom, \work{Matthew}, LXX, §§1--2.}

Aquinas's expressly mystical interpretation brings the coin's image into relation with the soul's image. The analogy has a limit that gives it force. An inscription can mark an artefact as part of an earthly economy. God does not merely stamp an already independent creature: the creature receives being from him. Rendering the person therefore cannot be reduced to transferring a religious portion of one's income while the rest of life becomes exempt. In Aquinas's exposition, even those spiritually raised above the world must seek peace in the world where they live with others. Belonging wholly to God can generate duties within shared life rather than dissolving them.\footnote{Aquinas, \work{Super Matthaeum}, 22, lectio 2, paragraph 87389, concluding literal and mystical explanations.}

The Gospel consequently offers a test rather than a ready-made modern constitution. Which real good does a claim serve, what is actually owed, and does compliance injure the allegiance due to God? Those questions require truthful discernment. They cannot be replaced by the assumption that every demand made by one's preferred ruler is lawful, or every demand made by an opponent is idolatrous. The narrated questioners' flattery warns that even true words about God can be recruited into a dishonest public argument.

\subsection{A people wider than an empire}
The responsorial Psalm supplies the widest horizon. All families of peoples are called to honour the Creator; their rulers and their achievements stand within creation. Augustine's new song is the song of a new love that builds a worldwide house. It crosses political and ethnic boundaries without making one earthly people the owner of divine glory. His warning to the preacher is essential here: the messenger who draws glory toward himself repeats within religion the appropriation that the Psalm exposes in idolatry.\footnote{Augustine, \work{Expositions}, XCV, §§2--6.}

Paul's thanksgiving shows what such a people looks like locally. The Thessalonians are not praised as a rival imperial administration. They are named as a community in God whose faith acts, love labours and hope endures. The independent apostolic reading thus contributes a different measure of public fruitfulness. Visible weakness and costly perseverance can reveal divine power; visible dominance need not. Chrysostom's account of their endurance prevents the universal Psalm from becoming only a splendid vision above the difficulties of actual neighbours.

The acclamation then gives that community an outward task. Holding the word of life and shining within the world are a form of witness, not a claim to own the world's light. Aquinas's participated light makes service possible without self-deification. The Christian who has something true to say must receive truth before bearing it. This is the religious counterpart to receiving a civic good without declaring its agent a saviour.

\subsection{From shelter to a salvation power cannot buy}
The Entrance contributes a necessary vulnerability. A person who must answer worldly claims may need courage, but courage does not make him invulnerable. The petition for hearing and shelter keeps fidelity from turning into a performance of self-sufficiency. Augustine's Christ-and-body reading places the assembly's confidence in a received mercy. The Collect directs the will toward sincere service, so that resistance to an excessive earthly claim does not become an excuse for a self-willed life.

At the altar, the Prayer over the Offerings adds purification. One can render a debt outwardly while inwardly seeking superiority. Even religious giving can become a means of control, as though possessing something to offer made the giver the master of the encounter. A petition for cleansing denies that conclusion without denying the gift's worth. Worship returns created goods to God while acknowledging that the worshipper needs more than the ability to give them.

The Psalm Communion option carries this movement through the failure of armies and horses to provide final deliverance. Augustine's nourishment of truth and Aquinas's sacramental food name a need beyond military security. The believer may seek ordinary protection and necessary food; he cannot make them guarantees against death. Fear and hope keep earthly prudence joined to dependence upon mercy. The prayer's bodily setting matters precisely because bodily strength is not enough.

If the Mark alternative is used, the contrast becomes Christ's saving self-gift. The Son does not save by making the disciples better competitors for domination. Bede first locates his bloodshed within a gift for the world, then names willing belief in his explanation of ``for many.'' Universal lordship does not eliminate the person's owed response. The one who receives liberation owes himself to the liberator. That connection between Bede's believing reception and Aquinas's return of the divine image belongs to the theological argument here; Bede's own passage explains Mark's ransom, not Matthew's coin.\footnote{Bede, \work{In Marci}, III, cols.236A--B; Aquinas, \work{Super Matthaeum}, 22.2.}

The Prayer after Communion finally orders present help toward eternal gifts. This releases earthly goods from an impossible task. They can be good without being final; they can be sought without being worshipped. A person can work honestly for the common life and still wait for a fulfilment that no ruler, institution or personal achievement can deliver. Such hope is not indifference to the world. It makes faithful action possible without demanding that history provide the last judgment upon its worth.

\subsection{The four senses}
\begin{fourSenses}
\item[Literal.] Cyrus is a real king used for Israel's restoration while remaining unknowing; Jesus answers an actual tribute trap; Paul thanks God for an actual community's active fidelity. The Psalms ask for protection and give the Creator glory. The prayers place sincere service, purification and sacramental benefit before God.
\item[Allegorical.] Cyrus's limited liberation anticipates the true Christ without transferring ignorance to him. Augustine's worldwide new song describes the gathered Church, and Aquinas's image reading places the whole person under God's claim. Where the Mark Communion option is used, the liberator is named through his self-gift and believing reception.
\item[Moral.] Render genuine duties truthfully, refuse to make impiety righteous by appealing to office, and return the whole self to God. The person who serves must also accept purification. Civic responsibility and worship meet in one life whose powers are received.
\item[Anagogical.] Equitable divine government and eternal gifts free hope from dependence on transient success. The broader Psalm contexts promise a fulfilment beyond death and incorruptible nourishment; their final horizon explains the worshipper's patience without promising present immunity from suffering.
\end{fourSenses}
